\section*{Problemas}

\subsection*{Especificación de conjuntos}

\begin{problema}
	Escriba en forma extensiva los siguientes conjuntos:
	\begin{enumerate}
		\item $A = \set{x \in \Z \mid -3 < x \leq 4}$
		\item $B = \set{x \in \N \mid x^{2} < 50}$
		\item $C = \set{x \in \Z \mid x \text{ es divisor de } 12}$
		\item $D = \set{x \in \R \mid x^{2} = 9}$
	\end{enumerate}
\end{problema}

\begin{problema}
	Escriba en forma intensiva los siguientes conjuntos:
	\begin{enumerate}
		\item $A = \set{2,4,6,8,10}$
		\item $B = \set{1,4,9,16,25}$
		\item $C = \set{\dots,-4,-2,0,2,4,\dots}$
	\end{enumerate}
\end{problema}

\begin{problema}
	\label{lip:exmp:1.4.a}
	Sea $A=\set{1,2,3,4},$ $B=\set{3,4,5,6,7},$ $C=\set{2,3,8,9}.$ Encuentre 
	\begin{enumerate}
		\item $A \cup B$ 
		\item $A \cap B$ 
		\item $A \cup C$ 
		\item $A \cap C$ 
		\item $B \cup C$ 
		\item $B \cap C$
	\end{enumerate}
\end{problema}

\subsection*{Subconjuntos}

\begin{problema}
	\label{thm:1.3}
	Demuestre que para cualesquiera dos conjuntos $A$ y $B,$ tenemos:
	$$
	A \cap B \subset A \subset A \cup B.
	$$
\end{problema}

\begin{problema}
	Determine cuáles de las siguientes afirmaciones son verdaderas:
	\begin{enumerate}
		\item $\emptyset \in \emptyset$
		\item $\emptyset \subset \emptyset$
		\item $\emptyset \in \set{\emptyset}$
		\item $\emptyset \subset \set{\emptyset}$
		\item $\set{1} \in \set{1,2}$
		\item $\set{1} \subset \set{1,2}$
	\end{enumerate}
\end{problema}

\begin{problema}
	Sea $A = \set{a, b, c}.$ Determine cuáles de las siguientes afirmaciones son verdaderas:
	\begin{enumerate}
		\item $a \in A$
		\item $\set{a} \in A$
		\item $\set{a} \subset A$
		\item $\emptyset \in A$
		\item $\emptyset \subset A$
		\item $\set{a,b} \subset A$
	\end{enumerate}
\end{problema}

\begin{problema}
	\label{thm:1.4}
	Demuestre que las siguientes proposiciones son equivalentes:
	\begin{enumerate}
		\item $\displaystyle A \subset B$
		\item $\displaystyle A \cap B = A$
		\item $\displaystyle A \cup B = B$
	\end{enumerate}
\end{problema}

\subsection*{Conjunto potencia}

\begin{problema}
	Calcule el conjunto potencia de:
	\begin{enumerate}
		\item $A = \set{1,2}$
		\item $B = \set{a,b,c}$
		\item $C = \emptyset$
	\end{enumerate}
\end{problema}

\begin{problema}
	Demuestre por inducción que si $\card{A} = n,$ entonces $\card{\mathcal{P}(A)} = 2^{n}.$
\end{problema}

\subsection*{Diagramas de Venn y operaciones}

\begin{problema}
	Sean $A = \set{1,2,3,4,5},$ $B = \set{3,4,5,6,7}$ y $C = \set{5,6,7,8,9}.$ Calcule:
	\begin{enumerate}
		\item $A \cup B \cup C$
		\item $A \cap B \cap C$
		\item $(A \cup B) \cap C$
		\item $A \cup (B \cap C)$
		\item $(A \setminus B) \cup (B \setminus A)$
	\end{enumerate}
\end{problema}

\begin{problema}
	En una clase de 40 estudiantes, 25 estudian matemáticas, 20 estudian física y 10 estudian ambas materias.
	\begin{enumerate}
		\item ?`Cuántos estudiantes estudian al menos una de las dos materias?
		\item ?`Cuántos estudiantes no estudian ninguna de las dos materias?
		\item ?`Cuántos estudian matemáticas pero no física?
	\end{enumerate}
\end{problema}

\begin{problema}
	Demuestre las siguientes identidades usando diagramas de Venn:
	\begin{enumerate}
		\item $A \cup (B \cap C) = (A \cup B) \cap (A \cup C)$
		\item $A \cap (B \cup C) = (A \cap B) \cup (A \cap C)$
		\item $A \setminus (B \cup C) = (A \setminus B) \cap (A \setminus C)$
	\end{enumerate}
\end{problema}

\begin{problema}
	Sea $U = \set{1,2,3,4,5,6,7,8,9,10},$ $A = \set{1,3,5,7,9},$ $B = \set{2,4,6,8,10},$ $C = \set{1,2,3,4,5}.$ Calcule:
	\begin{enumerate}
		\item $A^{C}$
		\item $A \cap B$
		\item $(A \cup C)^{C}$
		\item $A^{C} \cap C^{C}$
		\item Verifique que $(A \cup C)^{C} = A^{C} \cap C^{C}$ (Ley de De Morgan)
	\end{enumerate}
\end{problema}
